\section{Running websites with Flask}
\label{research:RunningWebsitesFlask}
To begin building my first prototype, I needed a way to display and interact with my website code directly on my laptop. In the past, I had only created very simple websites using static HTML files, which I would open directly in a browser. This approach was fine for basic structure and styling, but it quickly became limiting once I wanted to introduce dynamic content, user interaction, or database integration.

To find a more complete solution, I turned to web development communities like Stack Overflow and Reddit. These platforms gave me a clearer understanding of the range of tools available for creating full-stack applications. At this point, I didn't have much context on how professional web apps were typically built, so I took some time to learn about the different options and workflows.

Since Python is the language I'm most comfortable with, I focused my research on frameworks that use it. This led me to a few popular choices: Flask, Django and FastAPI

I decided to go with Flask for my prototype. It struck a good balance between simplicity and power. Using Flask, I could serve dynamic web pages that load data from a backend database — something static HTML simply can't do. This was especially important because I wanted my app to go beyond displaying fixed content. It needed to respond to user input, pull in components based on tags, and update content in real-time.

By starting with Flask, I gave myself a good foundation to develop a system that mimics how real-world web applications are built, while still staying within a language and environment I felt comfortable working in.

To learn the various systems within flask I spent some time reading through the documentation \cite{WelcomeFlaskFlask} and learning from some brief tutorials.

As seen in my prototypes, I used Flask to create a simple web server that can handle requests and serve HTML pages. This allows me to build a more interactive experience where users can submit data, which the server processes and responds to dynamically.
